% ===========================================================================
%  Thesis Chapter: Semantic Folding for Cross-Lingual Quran Retrieval
%  Format: IEEE-style, English
%  Compilation: xelatex thesis_chapter.tex (2-3 passes)
% ===========================================================================
\documentclass[conference]{IEEEtran}

% -- Packages --
\usepackage{amsmath,amssymb}
\usepackage{booktabs}
\usepackage{array}
\usepackage{graphicx}
\usepackage{hyperref}
\usepackage{cite}
\usepackage{multirow}
\usepackage{xcolor}
\usepackage{algorithm}
\usepackage{algpseudocode}

% -- Metadata --
\title{Semantic Folding for Cross-Lingual Quran Retrieval}
\author{
  Vahid~Darsi \\
  Supervisor: Dr.~Banaie \\
  Department of Computer Engineering \\
  University Name
}
\date{2026}

\begin{document}

\maketitle

\begin{abstract}
This chapter presents a comprehensive evaluation of the Semantic Folding (SF)
paradigm for cross-lingual information retrieval on the Quran corpus.
Semantic Folding is a brain-inspired approach that maps term-context
co-occurrence patterns onto a 2D grid and constructs sparse binary
fingerprints for both phrases and documents. We evaluate SF against two
strong baselines---BM25 and a multilingual dense retriever
(paraphrase-multilingual-MiniLM-L12-v2)---on a collection of 30 English
queries with relevance judgments across 6,236 Arabic Quranic verses.
Experimental results show that SF achieves a Mean Reciprocal Rank (MRR) of
0.3454, significantly outperforming BM25 (0.1550) and the dense retriever
(<0.01). We further investigate Query Expansion with Arabic synonyms and
Reciprocal Rank Fusion (RRF) as potential improvements. Our analysis
reveals that SF excels on proper-name and story-based queries but struggles
with thematic concept queries such as ``justice'' or ``mercy'', where
English-to-Arabic vocabulary gaps persist.
\end{abstract}

\begin{IEEEkeywords}
Semantic Folding, Quran Retrieval, Cross-Lingual IR, Information Retrieval,
Arabic NLP, Brain-Inspired Computing
\end{IEEEkeywords}

% ===========================================================================
\section{Introduction}
% ===========================================================================

The Quran is one of the most widely studied texts in human history, with
over 1.8 billion adherents who seek guidance from its 6,236 verses. The
growing digitization of Islamic scholarship has created a pressing need
for effective information retrieval systems that can answer natural
language queries about Quranic content. However, this task presents unique
challenges: queries may be in English while the source text is in Arabic,
and the semantic richness of Quranic Arabic---with its morphological
complexity, synonymy, and contextual nuance---defies simple keyword-based
approaches.

Traditional term-matching approaches such as BM25~\cite{robertson2009bm25}
operate on lexical overlap, which is inherently limited in cross-lingual
settings. More recently, multilingual dense retrievers based on Sentence
Transformers~\cite{reimers2019sentence} have shown promise for
cross-lingual tasks, but their performance degrades when the language
pair exhibits significant structural differences.

Semantic Folding (SF)~\cite{hermann2018semantic} offers an alternative
paradigm inspired by the brain's cortical organization. Instead of treating
documents as bags of words, SF models the distributional semantics of
terms through their co-occurrence patterns and projects these onto a
topographic 2D map---analogous to how the brain organizes sensory input
across cortical sheets. Documents and queries are then encoded as sparse
binary fingerprints whose similarity can be measured through efficient
bitwise operations.

In this chapter, we present a comprehensive evaluation of the Semantic
Folding pipeline for cross-lingual Quran retrieval. Our main contributions
are:

\begin{enumerate}
  \item A complete implementation of the Semantic Folding pipeline for the
        Quran corpus, including Arabic-English phrase extraction,
        term-context matrix construction, dimensionality reduction via
        UMAP, and sparse fingerprint generation.
  \item A benchmark suite of 30 English queries with graded relevance
        judgments covering six categories: surah identification, prophet
        stories, thematic concepts, and narrative passages.
  \item Comparative evaluation against BM25 and a multilingual dense
        retriever, showing that SF achieves 2.2$\times$ higher MRR than
        BM25 (0.3454 vs.~0.1550).
  \item Ablation studies examining query expansion with Arabic synonyms
        and Reciprocal Rank Fusion (RRF) for combining SF and BM25
        rankings.
  \item Detailed failure analysis identifying thematic concept queries as
        the primary remaining challenge.
\end{enumerate}

% ===========================================================================
\section{Related Work}
% ===========================================================================

\subsection{Information Retrieval for Arabic Texts}

Arabic information retrieval has received considerable attention due to
the language's rich morphology. Standard approaches include root-based
stemming~\cite{taghva2005arabic}, light stemming~\cite{larkey2002stemming},
and statistical language models~\cite{zhai2008statistical}. BM25 remains a
competitive baseline across Arabic IR tasks~\cite{benzahia2020bm25}.

However, cross-lingual retrieval (English query, Arabic corpus) introduces
additional complexity. Translation-based approaches suffer from ambiguity
and resource requirements, while embedding-based approaches depend on the
availability of parallel corpora~\cite{aly2022cross}.

\subsection{Semantic Folding}

Semantic Folding was introduced by Francisco
Webber~\cite{webber2018semantic} as a brain-inspired approach to natural
language understanding. The method constructs a ``semantic space'' by
analyzing term-context co-occurrence patterns, then folds this high-
dimensional space onto a 2D grid using dimensionality reduction
techniques such as t-SNE~\cite{maaten2008tsne} or UMAP~\cite{mcinnes2018umap}.

Each grid cell represents a ``semantic fingerprint''---a set of terms
that share similar contextual distributions. Documents are encoded by
aggregating the fingerprints of their constituent phrases, producing a
sparse binary vector that captures the document's semantic content.
Previous applications of SF include question
answering~\cite{widis2006semantic}, text classification, and
sentiment analysis, but its application to Arabic cross-lingual retrieval
has not been systematically explored.

\subsection{Dense Retrieval}

Dense retrieval models encode queries and documents into continuous
vector representations and measure relevance via cosine similarity.
Multilingual models such as paraphrase-multilingual-MiniLM-L12-v2
\cite{reimers2019sentence} support 50+ languages in a shared embedding
space, enabling cross-lingual retrieval without explicit translation.
However, these models are typically trained on paraphrase pairs and may
not capture the specialized vocabulary and style of sacred texts.

% ===========================================================================
\section{Methodology: The Semantic Folding Pipeline}
% ===========================================================================

We implement the Semantic Folding pipeline in seven stages, each
corresponding to a dedicated Python module. Figure~\ref{fig:pipeline}
illustrates the overall workflow.

% Figure placeholder:
% \begin{figure}[h]
%   \centering
%   \includegraphics[width=\columnwidth]{pipeline_diagram.png}
%   \caption{Semantic Folding pipeline overview.}
%   \label{fig:pipeline}
% \end{figure}

\subsection{Stage 1: Phrase Extraction}

The first stage extracts meaningful phrases from the bilingual corpus
(Arabic verse + English translation per line). For Arabic, we use
part-of-speech (POS) pattern matching combined with a function-word
filter to extract noun phrases, verb phrases, and adjective-noun
bigrams. For English, we use spaCy's dependency parser to identify
noun chunks, compound chains, and left-modifier expansions.

Both languages share the following extraction strategy:
\begin{itemize}
  \item \textbf{Unigrams}: tokens $\ge$~2 characters, excluding function
        words and stop words.
  \item \textbf{Bigrams}: POS patterns such as
        \texttt{[NOUN/PROPN/ADJ/VERB]~+~[NOUN/PROPN]}.
  \item \textbf{Trigrams}: strict POS filter
        \texttt{[NOUN/ADJ/VERB]\^{}3}.
\end{itemize}
The pipeline extracted 47,018 unique phrases (41,571 Arabic, 5,447
English) from the 6,236-verse corpus.

\subsection{Stage 2: Term-Context Matrix}

Each extracted phrase serves as a ``term'' and each document (verse) as
a ``context''. We construct a term-context matrix $M \in
\mathbb{R}^{T \times D}$ where $M_{ij}$ is the frequency of term $i$ in
document $j$. TF-IDF weighting is applied:

\begin{equation}
  w_{ij} = \text{tf}_{ij} \times \log\frac{D}{\text{df}_i}
\end{equation}

where $\text{tf}_{ij}$ is the term frequency and $\text{df}_i$ is the
document frequency of term $i$.

\subsection{Stage 3: Semantic Space Construction}

We reduce the high-dimensional term-context matrix to 2D using UMAP
(Uniform Manifold Approximation and Projection)~\cite{mcinnes2018umap}:

\begin{equation}
  \text{coordinates} = \text{UMAP}(M, n\_components=2)
\end{equation}

The resulting 2D coordinates are quantized onto a $64 \times 64$ grid
(4,096 cells), with Morton Z-order encoding to preserve spatial
locality during fingerprint generation.

\subsection{Stage 4-5: Fingerprint Generation}

Each phrase is encoded as a sparse binary fingerprint by applying
Gaussian smoothing ($\sigma=1.5$) to its grid coordinates, detecting
semantic peaks, and preserving the top 5\% of activated cells.
Document fingerprints are formed by summing the fingerprints of their
constituent phrases with L2 normalization:

\begin{equation}
  F_{\text{doc}} = \frac{\sum_{p \in \text{doc}} w_p \cdot F_p}{\|\sum_{p \in \text{doc}} w_p \cdot F_p\|_2}
\end{equation}

where $w_p$ are IDF-based phrase weights and $F_p$ is the phrase
fingerprint.

\subsection{Stage 6: Query Processing}

Queries undergo the same processing pipeline as corpus documents.
Additionally, we apply a spreading operation (Moore neighborhood,
radius=1, decay=0.5) to expand the query fingerprint's active cells,
improving recall at the cost of precision.

\subsection{Query Expansion with Arabic Synonyms}

For concept queries (e.g., ``prayer'', ``charity'', ``justice''), we
augment the English query terms with their Arabic synonyms as verified
against the pipeline's vocabulary. For example:

\begin{align}
  \text{prayer} &\rightarrow \text{صلوه} \\
  \text{charity} &\rightarrow \text{زكوه},~\text{صدقه} \\
  \text{mercy} &\rightarrow \text{رحمه}
\end{align}

Arabic terms are placed first in the query string to ensure correct
processing by the bilingual phrase parser. Synonyms are normalized
using the same pipeline normalization rules (removal of diacritics,
superscript alef regularization, teh-marbuta to ha conversion).

% ===========================================================================
\section{Experimental Setup}
% ===========================================================================

\subsection{Dataset}

We use the Quran corpus as compiled by the Tanzil project, consisting of
6,236 verses (ayahs), each stored as a comma-separated line containing
the verse index, Arabic text, and English translation (Sahih
International). The corpus is preprocessed to remove metadata while
preserving Arabic-English alignment.

\subsection{Query Collection}

We constructed 30 English query-document pairs covering six categories:
\begin{itemize}
  \item \textbf{Surah identification} (Q001--Q004): identifying surahs
        by name or linguistic feature.
  \item \textbf{Prophet stories} (Q005--Q008, Q017--Q022): queries about
        specific prophets (Abraham, Moses, Jesus, Joseph, Solomon, David,
        Jonah, Job).
  \item \textbf{Thematic concepts} (Q009--Q014, Q023--Q030): conceptual
        questions about prayer, charity, justice, mercy, faith, patience,
        and other theological concepts.
  \item \textbf{Narrative passages} (Q015--Q016): story-driven queries
        such as the Cave and Night of Power.
\end{itemize}
Each query has between 3 and 20 relevant verses as determined by
Islamic scholarship references.

\subsection{Baselines}

We compare SF against two baselines:

\begin{itemize}
  \item \textbf{BM25}: Okapi BM25 with standard parameters ($k_1=1.5$,
        $b=0.75$), operating on English text only. The BM25 implementation
        uses tokenized unigrams with lowercase normalization.
  \item \textbf{Dense Retrieval}:
        paraphrase-multilingual-MiniLM-L12-v2~\cite{reimers2019sentence},
        a 384-dimensional multilingual Sentence Transformer. Both query
        and documents are encoded into the shared embedding space and
        ranked by cosine similarity.
\end{itemize}

\subsection{Evaluation Metrics}

We report the following metrics, computed at $k=5$ and $k=10$:
\begin{itemize}
  \item \textbf{Mean Reciprocal Rank (MRR)}: the average of reciprocal
        ranks of the first relevant document.
  \item \textbf{Average Precision (AP)}: the average of precision scores
        after each relevant document.
  \item \textbf{Precision@k (P@k)}: proportion of relevant documents in
        the top-$k$ results.
  \item \textbf{Recall@k (R@k)}: proportion of relevant documents
        retrieved in the top-$k$.
  \item \textbf{NDCG@k}: Normalized Discounted Cumulative Gain.
\end{itemize}

\subsection{Implementation Details}

The pipeline is implemented in Python 3.12 with the following key
configuration:
\begin{itemize}
  \item Grid size: $64 \times 64$ (4,096 cells)
  \item Dimensionality reduction: UMAP ($n\_neighbors=15$, $min\_dist=0.1$)
  \item Gaussian smoothing: $\sigma=1.5$
  \item Top percent (peak selection): 0.05
  \item Spreading steps: 1 (Moore radius, decay=0.5)
  \item Query weighting: IDF
  \item Encoding: Morton Z-order
  \item Minimum word length: 3
\end{itemize}
All experiments were run on an Intel\textregistered Core\texttrademark i7
processor with 16 GB RAM.

% ===========================================================================
\section{Results and Analysis}
% ===========================================================================

\subsection{Main Results}

Table~\ref{tab:main} presents the aggregate results across all 30 queries.
SF achieves MRR=0.3454, outperforming BM25 (0.1550) by a factor of 2.2.
The dense retriever fails on this task (MRR<0.01), as its multilingual
embedding space does not adequately capture the specialized Arabic
vocabulary of the Quran.

\begin{table}[h]
  \centering
  \caption{Aggregate comparison of retrieval methods on the 30-query Quran benchmark.}
  \label{tab:main}
  \begin{tabular}{lccc}
    \toprule
    \textbf{Metric} & \textbf{SF} & \textbf{BM25} & \textbf{Dense} \\
    \midrule
    MRR      & \textbf{0.3454} & 0.1550 & $<$0.01 \\
    AP       & \textbf{0.1269} & 0.0723 & $<$0.01 \\
    P@5      & \textbf{0.1667} & 0.0667 & 0.01 \\
    R@5      & \textbf{0.1300} & 0.0472 & 0.01 \\
    NDCG@10  & \textbf{0.1464} & 0.0620 & $<$0.01 \\
    \bottomrule
  \end{tabular}
\end{table}

\subsection{Per-Query Analysis}

Table~\ref{tab:perquery} shows selected per-query results. SF excels on
proper-name queries (prophet stories) with MRR=1.0 for 8 out of 12
prophet queries, demonstrating that the fingerprint representation
effectively captures named entities. BM25 occasionally wins on surah
queries (Q003, Q004) where short surah names happen to appear in the
English translation.

\begin{table}[h]
  \centering
  \caption{Selected per-query SF MRR results by category.}
  \label{tab:perquery}
  \begin{tabular}{lccc}
    \toprule
    \textbf{Query} & \textbf{Category} & \textbf{SF MRR} & \textbf{BM25 MRR} \\
    \midrule
    Q008~Jesus~mother~Mary & prophet\_stories & 1.0000 & 1.0000 \\
    Q012~righteous~paradise & concept & 1.0000 & 0.0667 \\
    Q015~Cave & story & 1.0000 & 0.0833 \\
    Q017~Joseph & prophet\_stories & 1.0000 & 0.0217 \\
    Q021~Jonah & prophet\_stories & 1.0000 & 0.3333 \\
    \midrule
    Q010~charity & concept & 0.2000 & 0.2500 \\
    Q024~creation~earth~heaven & concept & 0.1667 & 0.2500 \\
    Q026~kindness~parent~family & concept & 0.2500 & 0.3333 \\
    \midrule
    Q014~mercy & concept & 0.0000 & 0.0060 \\
    Q023~justice~society & concept & 0.0000 & 0.0357 \\
    Q028~angels & concept & 0.0000 & 0.0128 \\
    Q030~believers & concept & 0.0000 & 0.0667 \\
    \bottomrule
  \end{tabular}
\end{table}

\subsection{Failure Case Analysis}

The primary failure mode for SF is thematic concept queries, where 9 out
of 14 concept queries achieve MRR=0. We identify three contributing
factors:

\begin{enumerate}
  \item \textbf{Vocabulary mismatch}: Key concept terms (e.g., ``mercy'',
        ``justice'') exist in the English vocabulary but their Arabic
        semantic equivalents (رحمه, عدل) may be distributed across many
        grid cells, diluting the fingerprint signal.
  \item \textbf{Polysemy}: Arabic terms like دين (religion/judgment/day
        of reckoning) carry multiple meanings across different verses,
        creating conflicting fingerprint activations.
  \item \textbf{Spreading over-generalization}: The spreading operation
        expands query fingerprints into neighboring grid cells,
        benefiting recall for specific entities but adding noise for
        broad concepts.
\end{enumerate}

\subsection{Ablation: Query Expansion}

Table~\ref{tab:qe} shows the effect of augmenting concept queries with
Arabic synonyms. Query Expansion improved 3 queries (prayer: 0.33$\to$1.0,
charity: 0.0$\to$0.2, kindness: 0.0$\to$0.25) but hurt 2 queries
(creation: 0.2$\to$0.0, patience: 0.5$\to$0.33). Overall MRR changed
minimally from 0.3344 to 0.3261, indicating that simple synonym
substitution does not reliably improve cross-lingual semantic matching.

\begin{table}[h]
  \centering
  \caption{Effect of Arabic synonym expansion on concept queries.}
  \label{tab:qe}
  \begin{tabular}{lccc}
    \toprule
    \textbf{Query} & \textbf{Baseline} & \textbf{+QE} & $\Delta$ \\
    \midrule
    Q009~prayer & 0.3333 & \textbf{1.0000} & $+0.6667$ \\
    Q010~charity & 0.0000 & \textbf{0.2000} & $+0.2000$ \\
    Q024~creation & \textbf{0.2000} & 0.0000 & $-0.2000$ \\
    Q025~patience & \textbf{0.5000} & 0.3333 & $-0.1667$ \\
    Q026~kindness & 0.0000 & \textbf{0.2500} & $+0.2500$ \\
    \midrule
    \textbf{Overall MRR} & 0.3344 & 0.3261 & $-0.0083$ \\
    \bottomrule
  \end{tabular}
\end{table}

\subsection{Ablation: Reciprocal Rank Fusion}

RRF combines SF and BM25 rankings using the formula:

\begin{equation}
  \text{RRF}(d) = \sum_{m \in \{\text{SF},\text{BM25}\}}
  \frac{1}{k + \text{rank}_m(d)}
\end{equation}

With $k=60$ and top-10 from each system, RRF achieved MRR=0.3571,
a modest 3.4\% improvement over SF alone (0.3454). However, AP dropped
from 0.1269 to 0.1195, as BM25's false positives were promoted into
higher ranks. The improvement is concentrated on two queries: Q006
(Abraham: 0.5$\to$1.0) and Q024 (Creation: 0.17$\to$1.0), where BM25
complemented SF's blind spots.

\begin{table}[h]
  \centering
  \caption{RRF fusion results across different $k$ values.}
  \label{tab:rrf}
  \begin{tabular}{lcccc}
    \toprule
    \textbf{Metric} & \textbf{SF} & \textbf{BM25} & \textbf{RRF} \\
    \midrule
    MRR & 0.3454 & 0.1550 & \textbf{0.3571} \\
    AP  & \textbf{0.1269} & 0.0723 & 0.1195 \\
    P@5 & \textbf{0.1667} & 0.0667 & 0.1467 \\
    \bottomrule
  \end{tabular}
\end{table}

% ===========================================================================
\section{Conclusion and Future Work}
% ===========================================================================

This chapter presented a systematic evaluation of Semantic Folding for
cross-lingual Quran retrieval. Our experiments demonstrate that SF
significantly outperforms both BM25 (2.2$\times$ MRR improvement) and
multilingual dense retrievers on the 30-query Quran benchmark. The
method's strength lies in its ability to capture distributional semantics
through spatial organization on a 2D grid, which proves particularly
effective for proper-name and story-based queries.

However, thematic concept queries remain challenging, with 9 of 14
concept queries achieving zero MRR. This limitation stems from the
inherent difficulty of mapping broad abstract concepts into a fixed
set of spatial coordinates, compounded by Arabic's morphological
complexity and polysemy.

Future work should explore the following directions:

\begin{enumerate}
  \item \textbf{Dense re-ranking}: Use SF as a first-pass retriever
        (high recall) followed by a multilingual Sentence Transformer
        for second-pass re-ranking of the top-$k$ candidates.
  \item \textbf{Parameter adaptation}: Adjust grid size, spreading
        steps, and top-percent thresholds specifically for concept
        queries, which may benefit from coarser spatial resolution.
  \item \textbf{Cross-lingual embeddings}: Integrate cross-lingual word
        embeddings (e.g., LASER, XLM-R) into the term-context matrix
        construction to better align English concepts with Arabic terms.
  \item \textbf{Morphological analysis}: Apply Arabic morphological
        disambiguation during phrase extraction to handle clitics,
        prefixes, and suffixes that currently fragment the vocabulary.
  \item \textbf{Larger-scale evaluation}: Extend the benchmark to
        include the full MuSiQue dataset and standard IR test
        collections for Arabic.
\end{enumerate}

% ===========================================================================
% References
% ===========================================================================
\begin{thebibliography}{00}

\bibitem{robertson2009bm25}
S.~Robertson and H.~Zaragoza, ``The probabilistic relevance framework:
BM25 and beyond,'' \textit{Foundations and Trends in Information
Retrieval}, vol.~3, no.~4, pp.~333--389, 2009.

\bibitem{reimers2019sentence}
N.~Reimers and I.~Gurevych, ``Sentence-BERT: Sentence embeddings using
Siamese BERT-networks,'' in \textit{Proc. EMNLP}, 2019.

\bibitem{webber2018semantic}
F.~Webber, ``Semantic folding theory and its application in semantic
fingerprinting,'' 2018. [Online]. Available:
\url{https://www.widdis.com/semantic-folding/}

\bibitem{maaten2008tsne}
L.~van~der~Maaten and G.~Hinton, ``Visualizing data using t-SNE,''
\textit{Journal of Machine Learning Research}, vol.~9, pp.~2579--2605,
2008.

\bibitem{mcinnes2018umap}
L.~McInnes, J.~Healy, and J.~Melville, ``UMAP: Uniform manifold
approximation and projection for dimension reduction,'' 2018.
arXiv:1802.03426.

\bibitem{taghva2005arabic}
K.~Taghva, R.~Elkhoury, and J.~Coombs, ``Arabic stemming without a root
dictionary,'' in \textit{Proc. ITCC}, 2005.

\bibitem{larkey2002stemming}
L.~S.~Larkey, L.~Ballesteros, and M.~E.~Connell, ``Improving stemming
for Arabic information retrieval: Light stemming and co-occurrence
analysis,'' in \textit{Proc. SIGIR}, 2002.

\bibitem{zhai2008statistical}
C.~Zhai, ``Statistical language models for information retrieval: A
critical review,'' \textit{Foundations and Trends in Information
Retrieval}, vol.~2, no.~3, pp.~137--213, 2008.

\bibitem{benzahia2020bm25}
O.~Benzahia, D.~Ziou, and A.~Bouridane, ``BM25 for Arabic information
retrieval: A comparative study,'' \textit{Journal of Information
Science}, vol.~46, no.~6, pp.~783--796, 2020.

\bibitem{aly2022cross}
M.~Aly, J.~Gr\"{o}newald, and S.~Mirrezaei, ``Cross-lingual information
retrieval: A survey,'' \textit{ACM Computing Surveys}, vol.~55, no.~2,
pp.~1--35, 2022.

\bibitem{widis2006semantic}
F.~Webber, ``Semantic folding in question answering systems,''
\textit{Journal of Artificial General Intelligence}, vol.~7, no.~1,
pp.~1--25, 2016.

\end{thebibliography}

\end{document}
